% Transcription — fonds Alexandre Grothendieck, shelfmark 156-7, batch 6 (pages 101-113).
% Established from the facsimile archives/batches/156-7/batch-06.pdf (a copy of
% 156-7.pdf fetched through the project's relay; no cover sheet in the batch,
% so PDF page k = folder page 100+k), per the skill transcribe-grothendieck. The
% folder is 113 pages in six batches; this is the sixth and last, thirteen
% pages. Batch 3 of this folder (pages 41-60, for the « Prat » axioms of page
% 51 and the example of page 53 that this batch cites) and batch 5 of the
% second draft (folder 156-6, GF VI) were read for notation; batches 4 and 5
% of this folder did not exist when this batch was written.
%
% Pass: Opus 5.5 (claude-opus-5-5), 2026-09-26 — first pass, unchecked against the pages by a human.
%
% Pages skipped: none. Pages summarised: none. All thirteen leaves are his
% mathematics, fast drafting in ink, and all are transcribed; page 113 stops
% a third of the way down, the end of the folder.
%
% His pagination 101-113 stands at the head of each leaf and coincides with
% the archivists' (noted on page 101). No date in his hand on these pages.
% Numbered runs, recorded where they stand: « Mag » axioms in three
% successive numberings — Mag Λ 1 ... Λ 6 / Mag 𝓜1-𝓜6 (formerly « Prat Λ »,
% p. 51), Mag 4'' and Mag 5'' (pp. 102-103), a second Mag 5'' (read so, the
% digit doubtful) and Mag 6'' (p. 106); a cancelled Mag 1-Mag 4 (p. 109) and
% the fair Mag 1-Mag 4 (p. 110); Mag fin (p. 113). Propositions 1, 2 (?) and
% an unnumbered Proposition (?) (pp. 103-104), a Lemme (p. 104), a Scholie
% (p. 106). Formula labels (*), (**), (***) (p. 106, the stars ill-formed).
% Atelier axioms At 1-At 3, At 4, a cancelled At 5, At 3 (spéc.) (a), (b)
% (pp. 111-112), At 𝓜 1-At 𝓜 4 (pp. 112-113). He cites At L1, L2, L3,
% At C (spéc.), At C L (spéc.), At C, At D. Page references: « p. 53 »
% (p. 101), « p. 51 » (pp. 102, 105), all to batch 3.
%
% What the batch is about. The atelier is re-founded on a new primitive.
% Page 101 proposes taking on a « magasin » (𝓜, ≤, ≪) the relation
% X |∘| Y ⇔ X° ∥ Y° (« intérieurement disjoints ») in place of ≬ or ∥, to
% give X° a sense even where the old definition makes it empty (the
% simplicial example: X° = ∅ unless X ∈ ℒ, « C'est triste mais c'est
% vrai »). Pages 102-105: ≬ defined from |∘| (X ≬ Y iff X' |∘| Y' for all
% strata outside the common closed part L), axioms Mag 4'' (distinct
% sub-strata of one X are |∘|), Mag 5'' (symmetric, antireflexive, stable
% under interior refinement), derivation of Mag 6, Proposition 1 (X ∥ Y iff
% all strata pairwise |∘|), Proposition 2 (At C (spéc.) as a condition C),
% a Lemme and the comparison with the definition of |∘| through interiors of
% shadows, under At L1-L3. Page 106: a Scholie — ≬-magasins satisfying
% L1-L3 and At C (spéc.) and |∘|-magasins satisfying the divisibility
% condition (***) correspond (p. 107, « équivalence de catégories »); p. 108
% explains why the |∘| direction is the richer one (it expresses an
% intuition that At L1 hides), and he « changes optics ». Pages 109-110:
% « Je reprends tout, une nouvelle fois ! » — a magasin (𝓜, ≤, ≪, |∘|) with
% Mag 1-Mag 4 (first set cancelled, then fair on p. 110), ≬ and ∥ defined,
% figures as ≤-closed pairwise compatible parts. Pages 111-113: ateliers
% (spéciaux) (𝔉, ≤, ≪; |∘|) with At 1-At 4 and At 3 (spéc.), then
% equivalently as 𝔉 ⊂ 𝔓(𝓜) with At 𝓜 1-4; ateliers « à figures finies »
% and Mag fin, and a closing remark that the atelier, not the set 𝔉 of all
% figures, is the centre of the structure.
%
% Relation to the earlier drafts. GF IV and GF VI take ≬ (and ∥) as the
% primitive relations; this batch replaces them by interior disjointness |∘|,
% from which ≬ and ∥ are defined (pp. 102, 110). Said in the notes where it
% happens rather than page by page; no « GF IV/VI, p. N » citation occurs on
% these pages.
%
% Reading hazards fixed once here, aligned with 156-6 and batch 3: 𝔉
% \mathfrak{F}, 𝓜 \mathcal{M}, ℒ \mathcal{L}; compatibility \between;
% disjointness \parallel; interior refinement \overset{\circ}{\ll}; « omb »,
% « Omb » in roman. New: his sign of two bars enclosing a small circle is set
% \mathrel{|{\circ}|}. The axioms he writes « Prot »/« Prat » are given
% « Prat » as batch 3 reads them on page 51. Circled (a), (b), (b') are set
% plain with a note. His underlining is \emph.
%
% Hard places, read only in part: the struck lines and margin of page 101;
% the slanted margins of pages 104, 109 and 111; page 108 (prose without
% formulas); the middle of page 111 (the description of 𝔉 in terms of
% (𝓜, ≤)); the end of page 107. Variances kept and flagged: « implique
% X |∘| Y » without primes in Mag 5'' (p. 103); « Mag 𝓜5 » for Mag 5''
% (p. 103); « Y'' = Ỹ' ∖ L' » and « Y' = Ỹ ∖ L » with = for ∈ (pp. 102,
% 104); « X'_L ≠ ∅ » for = ∅ (p. 104); the label Mag 5'' used twice (pp.
% 103, 106).

\documentclass[11pt,a4paper]{article}
\input{../preamble/grothendieck}

\folder{156-7}
\batch{6}
\pages{101}{113}
\dating{1986}
\watermark{Édition de démonstration}
\foldertitle{[Chapitre] VII. Analysis situs (troisième mouture) : notes manuscrites (23-26/06/1986)}

\begin{document}

\page{101}
\note{pagination de sa main 101 à 113 en tête des feuillets, qui coïncide avec celle des archivistes}

J'ai envie de prendre sur un « magasin » $(\mathcal{M}, \leq, \ll \ldots)$ comme relations primitives, en plus de $\leq$ et de $\ll$, et au lieu de $\between$ ou de $\parallel$, la \struck{nouv} relation (heuristiquement)
\[
  X \mathrel{|{\circ}|} Y \Longleftrightarrow X^{\circ} \parallel Y^{\circ}
\]
et ceci de façon à arriver à donner un sens raisonnable à $X^{\circ}$, même dans des cas où au sens pris jusqu'à présent \uncertain{celle-ci} serait vide. \uncertain{Cf.} p.~53 pour un exemple\note{le mot « magasin », entre guillemets, est nouveau dans ce lot ; le signe $\mathrel{|{\circ}|}$, deux traits verticaux enserrant un petit rond, se lit (page 109) « $X$ et $Y$ sont intérieurement disjoints ». Le renvoi « p. 53 » est au lot 3 de ce dossier ; devant, une grande lettre (un C ?) et un signe bouclé, lus « Cf. » sous réserve} : \struck{pour} \add{pour} une \uncertain{maquette} simpliciale, plus gén.\ un ensemble ordonné $(\mathcal{M}, \leq)$ \uncertain{où on} prend $\ll \;=\; \leq$, $X \between Y$ toujours satisfait) \struck{dans \ill{} pour tout $X$ \ill{} \ill{} un minimal \ill{}} \uncertain{si} $X$ \uncertain{est} un minimal \ill{} on a $\mathrm{omb}(X)^{\circ} = \emptyset$, \struck{\ill{} $X \smallsetminus \{X\}$}.\note{deux lignes et demie barrées d'un long trait, lues par fragments ; l'exposant de « $\mathrm{omb}(X)^{\circ}$ » ressemble à un $d$, lu comme le rond de l'intérieur, douteux. Un « V » sous le premier mot biffé de la ligne appelle l'insertion « pour »}

P.\ ex.\ pour une \uncertain{maquette} simpliciale $(\mathcal{L}, \mathcal{M} \subset \mathfrak{P}_f(\mathcal{L}))$ les supports de $(\mathcal{M}, \leq, \ll \;=\; \leq, \between \;=\; \text{toujours satisfait})$, pour la relation $X \parallel Y \Longleftrightarrow \widetilde{X} \cap \widetilde{Y} = \emptyset$ i.e.\ $X$ et $Y$ non voisins, s'identifient avec \uncertain{les} parties quelconques de $S$
\[
  \Sigma_{\mathcal{M}} \simeq \mathfrak{P}(S)
\]
\marginal{$\mathcal{M}$ partie \ill{} de $\mathfrak{P}_f(\mathcal{L})$ \ill{} $(\mathcal{L})$ pour \ill{} contenant les $\{x\}$}\note{note écrite en biais dans la marge gauche, en face de l'exemple, lue en partie ; « $(\mathcal{L}, \mathcal{M}$ » : les deux lettres sont repassées d'une encre épaisse}
le support du \struck{$X \in$} « simplexe » $X$ est $X$ \uncertain{lui-même}. Avec la définition plausible
\[
  X^{\circ} =: \mathrm{supp}\, X \cap \mathrm{cosupp}\, \partial X
\]
on trouve $X^{\circ} = \emptyset$ sauf si $X$ \struck{est} $\in \mathcal{L}$.

C'est triste mais c'est vrai.

\page{102}
Je vais prendre comme relations primitives dans $\mathcal{M}$
\[
  \leq, \quad \ll, \quad \mathrel{|{\circ}|}
\]
Avec ça, on va définir $\between$ ainsi (version \struck{\ill{}} « magasin spécial », on reviendra \uncertain{bientôt} sur la version « magasin polyédral ») ; pour $X, Y \in \mathcal{M}$, posant $L = \widetilde{X} \cap \widetilde{Y}$ \add{plus grande partie fermée commune}
\[
  X \between Y \Longleftrightarrow \forall\, X' \in \widetilde{X} \smallsetminus L,\ Y' \in \widetilde{Y} \smallsetminus L,\quad X' \mathrel{|{\circ}|} Y'
\]
\marginal{cette définition ne fait pas intervenir $\ll$}\note{note dans la marge gauche, en face de la définition ; « plus grande partie fermée commune » est écrit au-dessus de la formule}

Nous allons d'abord poser les axiomes Mag $\Lambda$ 1\ldots\add{\uncertain{$\Lambda$ 6}} (anciennement Prat $\Lambda$ 1 \ill{}\ldots) de la p.~51, avec $\Lambda = \mathcal{M}$.\note{« $\Lambda$ 6 » est écrit au-dessus de la ligne avec un signe d'insertion ; après « Prat $\Lambda$ 1 », une petite lettre bouclée (un $\xi$ ?) n'est pas lue ; les axiomes « Prat » (préatelier) et la page 51 sont au lot 3 de ce dossier ; le mot se lit aussi « Prot », il est donné comme au lot 3}
\begin{itemize}
  \item[] Mag $\mathcal{M}$1, Mag $\mathcal{M}$3, Mag $\mathcal{M}$4 en variantes \add{devenant Mag 1'', 2'', 3''}
  \item[] Mag $\mathcal{M}$2 tombe (car tautologique si $\Lambda = \mathcal{M}$)
  \item[] Mag $\mathcal{M}$5 \struck{est tautologique} va résulter des nouveaux axiomes (remplaçant At L3)
\end{itemize}
\note{au-dessus de « Mag $\mathcal{M}$4 », un « $\mathcal{M}\Lambda 2$ » relié par un trait courbe à la ligne suivante ; « devenant Mag 1'', 2'', 3'' » est écrit au-dessus de la ligne, à droite}

\emph{Mag} 4'' $\forall\, X \in \mathcal{M}$, et $Y, Z \leq X$, \add{$Y \neq Z$,} on a $Y \mathrel{|{\circ}|} Z$.\note{après le premier $Y$ de la conclusion, un signe lourdement biffé}

\begin{tikzcd}[column sep=small, row sep=small]
  & \widetilde{X} \arrow[dl, no head] \arrow[dr, no head] & & \widetilde{Y} \arrow[dl, no head] \arrow[dr, no head] & \\
  \widetilde{X}' \arrow[dr, no head] & & L \arrow[dl, no head] \arrow[dr, no head] & & \widetilde{Y}' \arrow[dl, no head] \\
  & X'_L \arrow[dr, no head] & & Y'_L \arrow[dl, no head] & \\
  & & L' & &
\end{tikzcd}

\drawing{schéma en losange à gauche de la démonstration ; à gauche, « $X'' \in$ » pointe vers $\widetilde{X}'$, à droite « $\ni Y''$ » vers $\widetilde{Y}'$}

Supposons $X \between Y$, soient $X' \leq X$, $Y' \leq Y$, prouvons $X' \between Y'$. Introduisons \struck{$\between$} $L = \widetilde{X} \cap \widetilde{Y}$, $L' = \widetilde{X}' \cap \widetilde{Y}'$, $X'_L = \widetilde{X}' \cap L = \widetilde{X}' \cap \widetilde{Y}$, $Y'_L = \widetilde{Y}' \cap L = \widetilde{Y}' \cap \widetilde{X}$, donc $L' = X'_L \cap Y'_L$. Soit $X'' \in \widetilde{X}' \smallsetminus L'$, $Y'' = \widetilde{Y}' \smallsetminus L'$, prouvons $X'' \mathrel{|{\circ}|} Y''$. Si $X'', Y'' \notin L$, cela \struck{résulte} de l'hypothèse que $X \between Y$. Si $X'' \in L$ i.e.\ $X'' \in \widetilde{Y}$ ($Y''$ \uncertain{non} \ill{}) alors $X'', Y'' \in \widetilde{Y}$ et $X'' \neq Y''$, donc par Mag 4'' on gagne encore.\note{« $Y'' = \widetilde{Y}' \smallsetminus L'$ » : un $=$ sur la page, pour $\in$. Le $\widetilde{Y}$ de « $X'' \in \widetilde{Y}$ » porte un indice biffé ; les mots entre parenthèses qui suivent ne sont pas lus}

Il faut encore \uncertain{adjoindre} Mag $\mathcal{M}$6 (ex Prat $\Lambda$6, pour $\Lambda = \mathcal{M}$).

Cor.\ de \emph{Mag} 4'' Si $X, Y$ avec $X \between Y$, $X \neq Y$, alors $X \mathrel{|{\circ}|} Y$.\note{le signe de « $X \between Y$ » est vite tracé et pourrait être un $\geq$ barré ; lu comme la compatibilité, douteux}

\page{103}
On voudra aussi la \uncertain{condition} qui suit, similaire à ex-Mag $\Lambda$5 :

\emph{Mag} 5'' \struck{La} La relation $X \mathrel{|{\circ}|} Y$ est symétrique et antiréflexive, i.e.\ elle implique $Y \mathrel{|{\circ}|} X$ et $X \neq Y$. De plus, si $X \mathrel{|{\circ}|} Y$ et $X' \mathrel{\overset{\circ}{\ll}} X$, $Y' \mathrel{\overset{\circ}{\ll}} Y$ implique $X \mathrel{|{\circ}|} Y$.\note{« implique $X \mathrel{|{\circ}|} Y$ » est sur la page, sans les primes qu'on attend ($X' \mathrel{|{\circ}|} Y'$), et la phrase mêle « si » et « implique »}

Ceci va impliquer déjà Mag 6, i.e.\ si $X, Y, \add{Z} \in \mathcal{M}$, $X \between Y$, $Z \ll X, Y$, alors $\exists\, Z' \in \mathcal{M}$ avec
\[
  Z \ll Z' \leq X, Y .
\]
\note{le « 6 » de « Mag 6 » est surchargé ; le $Z$ est inséré au-dessus de la ligne ; le premier $\ll$ de « $Z \ll X, Y$ » est écrit sur un autre signe}
En effet, par Mag $\mathcal{M}$3, \struck{OPS $Z \mathrel{\overset{\circ}{\ll}} X$, $Z \mathrel{\overset{\circ}{\ll}} Y$} $\exists\, X' \in \widetilde{X}$, $Y' \in \widetilde{Y}$ avec $Z \mathrel{\overset{\circ}{\ll}} X'$, $Z \mathrel{\overset{\circ}{\ll}} Y'$. Par \emph{Mag $\mathcal{M}$5} \add{qu'on vient de vérifier}, on aura $X' \between Y'$, et par le corollaire de Mag 4'', on conclut $X' = Y'$, i.e.\ le $Z'$ cherché !\note{« Mag $\mathcal{M}$5 » est souligné et relié par un trait à « qu'on vient de vérifier » ; on attend plutôt Mag 5'', énoncé en tête de page. Le signe de « $X' \between Y'$ » est mal formé}

Donc on sait qu'\struck{\ill{}} avec $(\leq, \ll, \between)$ on a un atelier, \struck{formé} dans \uncertain{lequel} \ill{} une relation $\parallel$

\emph{Proposition} 1 Soient $X, Y \in \mathcal{M}$. Alors
\[
  X \parallel Y \Longleftrightarrow \forall\, X' \in \widetilde{X},\ Y' \in \widetilde{Y}, \text{ on a } X' \mathrel{|{\circ}|} Y'
\]
En effet, $X \parallel Y$ signifie \add{$X \between Y$ et} \struck{que} $L = \widetilde{X} \cap \widetilde{Y} = \emptyset$. \struck{et $X \between Y$} Donc la prop.\ est tautologique.

Vérifions At C (spéc.), i.e.

\emph{Proposition} 2 \add{(?)} Soient $X, Y \in \mathcal{M}$, $L = \widetilde{X} \cap \widetilde{Y}$. \struck{Soi} Alors $X \between Y \Longleftrightarrow$ \add{(condition C)} $\forall\, X' \ll X$, $Y' \ll Y$, tels que $X'_L \overset{\mathrm{df}}{=} \widetilde{X}' \cap \mathrm{Omb}(L)$ et $Y'_L \overset{\mathrm{df}}{=} \widetilde{Y}' \cap \mathrm{Omb}(L)$ soient vides, on a $X' \parallel Y'$.\note{« (condition C) » est écrit sous « Soient $X, Y$ » et relié au second membre de l'équivalence ; avant $X'_L$, un premier essai « $X'_L = \widetilde{X}' \cap$ » est biffé ; les indices de « $X' \parallel Y'$ » sont lourdement biffés. « At C (spéc.) » renvoie à un axiome des lots précédents}

$\Longrightarrow$ semble toujours vrai, et $\Longleftarrow$ moyennant un axiome \struck{\ill{}} supplémentaire genre At L1 (axiome de divisibilité infinie). Voyons !

\page{104}
Supposons $X \between Y$, et soient $X' \ll X$, $Y' \ll Y$ avec
\[
  X'_L = Y'_L = \emptyset
\]
prouvons que $X' \parallel Y'$, i.e.\ \add{(prop.\ 1)} que \struck{pour} $\forall\, X'' \leq X'$, $Y'' \leq Y'$, on a $X'' \mathrel{|{\circ}|} Y''$. Comme \struck{$X''$} $X''_L \leq X'_L = \emptyset$ et \struck{\ill{}} $Y''_L = \emptyset$, OPS $X'' = X'$, $Y'' = Y'$. Soient $X_1$, \add{$Y_1$} tels que \struck{$X' \mathrel{\overset{\circ}{\ll}} X_1 \leq X$,}
\[
  X' \mathrel{\overset{\circ}{\ll}} X_1 \leq X, \qquad Y' \mathrel{\overset{\circ}{\ll}} Y_1 \leq Y
\]
\note{en tête de cette ligne, à gauche, un signe isolé (un $\exists$ ?)}
On a $X_1 \notin L$ (sinon on aurait $X' \in L$ et $X'_L \neq \emptyset$) et de même $Y_1 \notin L$, donc $X_1 \mathrel{|{\circ}|} Y_1$, par hyp.\ $X \between Y$. Donc par Mag 5'', on conclut $X' \mathrel{|{\circ}|} Y'$, cqfd.\note{« $X_1 \mathrel{|{\circ}|} Y_1$ » : le signe est serré, comme un H ; lu comme $\mathrel{|{\circ}|}$, que demande la définition de $\between$ (page 102). Dans « $X' \in L$ », le $\in$ est surchargé}

Supposons d'autre part l'hypothèse dite sur $X, Y$, en termes des $X', Y'$, satisfaite, prouvons que $X \between Y$. Soit donc $X' \in \widetilde{X} \smallsetminus L$ \struck{\ill{}}, $Y' = \widetilde{Y} \smallsetminus L$, prouvons $X' \mathrel{|{\circ}|} Y'$.\note{« $Y' = \widetilde{Y} \smallsetminus L$ » : un $=$ pour $\in$, comme page 102}

\struck{Ici on tombe \ill{} qu'il \ill{} soit} \uncertain{Si} \ill{}, \ill{}

\marginal{Si on suppose At L2 (?) \add{\ill{}} la condition C (\ill{}) \ill{} plus \ill{} : \struck{$\mathrm{omb}(X)$} $\forall\, x \in \mathrm{omb}(X)^{\circ}$, $y \in \mathrm{omb}(Y)^{\circ}$, \struck{\ill{}} on a $x \parallel y$ — i.e.\ \ill{} \ill{} pour tout $x \in$ \ill{}, $y \in$ \ill{}}\note{longue note écrite en biais dans la marge gauche, de la ligne « prouvons que $X \between Y$ » au lemme, lue par fragments ; une flèche la relie au texte}

\emph{Lemme} La condition C sur $(X, Y)$ implique que pour tous $X', Y'$ comme dans l'énoncé de cette condition, et pour $x \in \mathrm{omb}(X')^{\circ}$, $y \in \mathrm{omb}(Y')^{\circ}$, on a $x \parallel y$.\note{le lemme est marqué d'un trait vertical dans la marge}

En effet, on a bien sûr \struck{$x$} $x_L = \emptyset$ \struck{$y_L = \emptyset$}, car sinon on aurait $x_L = x$ donc $x \ll L$ \struck{donc} donc $\mathrm{omb}(X') \cap \mathrm{omb}(L) = \mathrm{omb}(X'_L) \neq \emptyset$, contrairement à \add{$X'_L = \emptyset$} la condition. Et de même $y_L = \emptyset$, donc \ill{}.\note{« contrairement à » est suivi, au-dessus de la ligne, d'un « $X'_L \neq \emptyset$ » (on attend $= \emptyset$) et, sur la ligne, de « la condition »}

Donc inversement, la condition At C (spéc.) signifie \struck{la condition} \add{celle de C} \struck{que si $X, Y \in \mathcal{M}$, $\Longrightarrow$} \struck{on aurait} soit (ex-$X'$, $Y'$) \struck{L'axiome de}

\emph{Proposition} (?) Soient $X, Y \in \mathcal{M}$. Alors
\[
  X \mathrel{|{\circ}|} Y \Longleftrightarrow \forall\, x \in \mathrm{omb}(X)^{\circ},\ y \in \mathrm{omb}(Y)^{\circ}, \text{ on a } x \parallel y .
\]

\page{105}
Bien sûr on a $\Longrightarrow$ à cause de Mag 5''. Il faut voir si on a $\Longleftarrow$. Si on part d'un magasin style $(\leq, \ll, \between)$, décrit par Mag 1 -- Mag 6 (ex Prat 1 -- Prat 6) ou Mag $\mathcal{M}$1 -- Mag $\mathcal{M}$6 (\ill{} la p.~51), et si on suppose \ill{} L1, L2, L3, alors on \struck{\uncertain{pourra}} \emph{définir} raisonnablement
\[
  X \mathrel{|{\circ}|} Y \Longleftrightarrow \forall\, x \in \mathrm{omb}(X)^{\circ},\ y \in \mathrm{omb}(Y)^{\circ}, \text{ on a } x \parallel y,
\]
et Mag 4'' et Mag 5'' sont vérifiés (Mag 4'' \uncertain{en raison} de At L3, et Mag 5'' \struck{pr}\uncertain{provenant} de $X' \mathrel{\overset{\circ}{\ll}} X \Rightarrow \mathrm{omb}(X')^{\circ} \subset \mathrm{omb}(X)^{\circ}$).\note{« Mag 1 -- Mag 6 » : le « 1 » est au-dessus d'un trait ; « Prat 6 » est écrit sur un autre chiffre ; « Mag $\mathcal{M}$1 -- Mag $\mathcal{M}$6 » est souligné d'un trait courbe qui le relie à « L1, L2, L3 ». Le mot qui précède « définir » (souligné) est surchargé ou biffé, lu « pourra » sous réserve}
Si on regarde alors la relation $\between'$ définie par cette relation $\mathrel{|{\circ}|}$ comme ci-dessus, alors on sait que cette relation et la relation $\between$ \ill{} ssi $X$ satisfait la condition At C (spéc.), équivalente \struck{à} At C L (spéc.) en effet, moyennant L2, i.e.\ à la condition que pour $X, Y \in \mathcal{M}$, $L = \widetilde{X} \cap \widetilde{Y}$
\[
  X \between Y \Longleftrightarrow \underbrace{\forall\, x \in \mathrm{omb}(X) \smallsetminus \mathrm{omb}(L),\ y \in \mathrm{omb}(Y) \smallsetminus \mathrm{omb}(L), \text{ on a } x \parallel y}_{X \between' Y}
\]
\note{« At C (spéc.) » : le « At » est écrit sur un autre mot ; entre « L2, » et « At C L », un signe biffé. « $X \between' Y$ » est écrit sous l'accolade du second membre}

En effet, la condition du second membre \ill{} équivaut à $X \between' Y$, car
\[
  \mathrm{omb}(X) \smallsetminus \mathrm{omb}(L) = \bigcup_{X' \in \widetilde{X} \smallsetminus L} \mathrm{omb}(X')^{\circ}, \qquad
  \mathrm{omb}(Y) \smallsetminus \mathrm{omb}(L) = \bigcup_{Y' \in \widetilde{Y} \smallsetminus L} \mathrm{omb}(Y')^{\circ}
\]
donc la condition \add{du} second membre s'écrit aussi :
\[
  \forall\, X' \in \widetilde{X} \smallsetminus L,\ Y' \in \widetilde{Y} \smallsetminus L, \text{ on a } \underbrace{\bigl[\forall\, x \in \mathrm{omb}(X')^{\circ},\ y \in \mathrm{omb}(Y')^{\circ} \text{ on a } x \parallel y\bigr]}_{X' \mathrel{|{\circ}|} Y'}
\]
i.e.\ on a bien $X \between' Y$.

\page{106}
En résumé :

\emph{Scholie} sur la comparaison des points de vue « Magasin $\leq, \ll, \between$ » et « Magasin $\leq, \ll, \mathrel{|{\circ}|}$ ».

Partant d'un $\between$-magasin, satisfaisant L1, L2, L3, celui-ci définit un $\mathrel{|{\circ}|}$-magasin, en définissant
\[
  (*) \qquad X \mathrel{|{\circ}|} Y \Longleftrightarrow \forall\, x \in \mathrm{omb}(X)^{\circ},\ y \in \mathrm{omb}(Y)^{\circ}, \text{ on a } x \parallel y
\]
D'autre part, tout $\mathrel{|{\circ}|}$-magasin (\uncertain{certain}) définit un $\between$-magasin, où \add{\uncertain{donc}} la relation $\parallel$ sera équivalente à
\[
  (**) \qquad X \parallel Y \Longleftrightarrow \forall\, X' \in \widetilde{X},\ Y' \in \widetilde{Y}, \text{ on a } X' \mathrel{|{\circ}|} Y' .
\]
\add{(il semble que)}
\marginal{$(**)$ $X \between Y \Longleftarrow \forall\, X' \in \widetilde{X} \smallsetminus \widetilde{X} \cap \widetilde{Y}$, $Y'' \in \widetilde{Y} \smallsetminus \widetilde{X} \cap \widetilde{Y}$, on a $X' \mathrel{|{\circ}|} Y'$}\note{note écrite verticalement dans la marge gauche, reliée par une accolade à la ligne « $\between$-magasin » et à $(**)$ ; « (il semble que) » est écrit sous le second membre de $(**)$ et renvoyé à la phrase suivante. Les étoiles des étiquettes $(*)$, $(**)$ sont mal formées}

Mais on s'intéresse (la relation $\mathrel{|{\circ}|}$ un peu au-delà de la définition raisonnable en termes des \uncertain{structures} $\leq, \ll, \between$) à ce qu'elles donnent raisonnablement. Pour qu'elle le puisse, il faut supposer qu'elle satisfait la condition supplémentaire

\marginal{implique At L1}\note{en biais dans la marge gauche, en face de l'axiome}
\emph{Mag} 5'' (du type d'un axiome de divisibilité) Si $X, Y \in \mathcal{M}$, alors
\[
  (***) \qquad X \mathrel{|{\circ}|} Y \Longleftrightarrow \forall\, x \in \mathrm{omb}(X)^{\circ},\ y \in \mathrm{omb}(Y)^{\circ}, \text{ on a } x \mathrel{|{\circ}|} y
\]
(\struck{\ill{}} en tous cas, pour $x, y \in \mathcal{L}$, $x \mathrel{|{\circ}|} y \Leftrightarrow x \parallel y$) --- i.e.\ que la relation $X \mathrel{|{\circ}|} Y$ provienne d'une relation (qu'on peut aussi noter $\parallel$) \emph{sur $\mathcal{L}$}. Ceci n'est possible \emph{que} si l'atelier satisfait la condition de divisibilité, i.e.\ Mag 6'' \ill{}\note{l'étiquette de l'axiome est lue « Mag 5'' », déjà prise page 103 ; le « 5 » est mal formé et pourrait être un autre chiffre. L'étiquette de la formule porte trois ou quatre étoiles. Le dernier mot de la page n'est pas lu ; la phrase continue en tête de la page 107}

\page{107}
pour
\[
  \mathrm{omb}(X)^{\circ} \neq \emptyset \qquad \forall\, X \in \mathcal{M},
\]
car si on avait $\mathrm{omb}(X)^{\circ} = \emptyset$, on aurait $X \mathrel{|{\circ}|} X$, ce qui est exclu par Mag 5'' (antiréflexivité). De plus, elle implique aussi At L2 (via $(***)$) et At L3 (via \struck{Mag} Cor.\ de Mag 4'').\note{l'étiquette « $(***)$ » renvoie à la formule de la page 106, où elle porte trois ou quatre étoiles}

Et l'axiome At C (spéc.) semble automatique (alors qu'il ne l'est peut-être pas, et qu'il est bien, dans le cas où \ill{} \ill{}, supposé par $(***)$).

Inversement, si on part d'un $\between$-magasin, satisfaisant At L1, L2, L3 (qui permettent y définir raisonnablement $\mathrm{|{\circ}|}$), la structure $\between$ \ill{} \struck{\ill{}} \add{$\between'$} définie par $\mathrel{|{\circ}|}$, ssi autre que elle, \ill{} $\mathcal{M}$ satisfait At C (spéc.).\note{la fin de cette phrase est mal lue : on attend que $\between$ et la relation $\between'$ définie par $\mathrel{|{\circ}|}$ coïncident ssi $\mathcal{M}$ satisfait At C (spéc.), comme à la page 105}

Donc la \uncertain{donnée} \ill{} commune des deux notions est, visiblement \ill{} par l'équivalence de catégories :

\begin{tikzcd}[row sep=small]
  \between\text{-magasins satisfaisant At L1, L2, L3 et At C (spéc.)} \arrow[d, Rightarrow] \\
  \mathrel{|{\circ}|}\text{-magasins satisfaisant } (***) \text{ (compatibilisme)}
\end{tikzcd}

\note{la flèche verticale double entre les deux lignes n'a de pointe visible que vers le bas (le haut est un crochet ouvert) ; elle est rendue telle quoique le texte annonce une équivalence. « (dispositions) », entre parenthèses et souligné, est écrit à droite de la seconde ligne et relié à la suite}

Dans les $\between$-magasins et les $\mathrel{|{\circ}|}$-magasins apparaissent comme deux \struck{types de} généralisations d'une même espèce de structure. Celle-ci paraît trop \uncertain{restrictive} \ill{}, en un sens, par la seule présence de At L1.

\page{108}
On dispose donc de deux généralisations possibles, en direction d'une topologie sans axiome de divisibilité. Mais c'est la direction $\mathrel{|{\circ}|}$ qui me semble la plus \uncertain{riche}. En effet, elle a toujours une structure de $\between$-magasin correspondante, mais, en \emph{plus}, une notion $\mathrel{|{\circ}|}$ de disjonction \uncertain{intérieure}, qui exprime une \uncertain{intuition} géométrique \uncertain{inexprimable}, laquelle \uncertain{échappe} (vu l'absence de At L1) au seul $\between$-\struck{point} \ill{}. L'axiome At C (spéc.), dans le cas \uncertain{actuel}, \uncertain{apparaît} artificiel --- il exprime \ill{} l'absence de \ill{} celui qui \ill{} dispose d'une notion de « disjonction » \ill{} les uns des \uncertain{intérieurs} « strates » \ill{} $X^{\circ}$.\note{page d'écriture rapide et sans formules ; seuls les mots offerts sont lus, souvent sous réserve}

Donc finalement, je change d'optique \ill{} celles d'axiomes de \ill{} \ill{} At C (spéc.), At C, At D\ldots\ tombent \ill{} \ill{} \uncertain{flèches}. Je reprends\ldots

\page{109}
tout, une nouvelle fois !

Un \emph{magasin} (de multistrates) est la donnée d'un ensemble $\mathcal{M}$, muni de trois relations
\[
  \leq, \quad \ll, \quad \mathrel{|{\circ}|}
\]
avec la terminologie introduite par ailleurs. La relation $X \mathrel{|{\circ}|} Y$ se lit « $X$ et $Y$ sont intérieurement disjoints », on pourra aussi l'écrire plus loin $X^{\circ} \parallel Y^{\circ}$ (quand on disposera des notations : $X^{\circ}$, $Y^{\circ}$\ldots).

Les axiomes sont les suivants\note{toute la suite de la page (Mag 1 à Mag 4 et la note de la marge) est barrée de longs traits obliques ; elle est transcrite en entier}

\struck{\emph{Mag 1} $X \leq Y$ et $X \ll Y$ sont deux relations d'ordre, et $X \leq Y \Longrightarrow X \ll Y$}

\struck{\emph{Mag 2} $\forall\, X \in \mathcal{M}$, $Z \in \mathcal{M}$ avec $Z \ll X$, $\exists\, X'$ avec $Z \mathrel{\overset{\circ}{\ll}} X' \leq X$, où on écrit $Z \mathrel{\overset{\circ}{\ll}} Y \overset{\mathrm{def}}{\Longleftrightarrow} Z$ est élément minimal de $\widetilde{Y}^{Z} = \lbrace Y' \in \widetilde{Y} \mid Z \ll Y' \rbrace$}\note{après « $\forall\, X \in \mathcal{M}$, » un signe lourdement biffé ; dans « $\widetilde{Y}^{Z}$ » le $Y$ est écrit sur un $X$}

\struck{\emph{Mag 3} Soient $Y \ll X$ et $Z$ tels que $Z \ll Y, X'$ \ill{}, alors $\exists\, Y'$ avec $Z \ll Y' \ll X'$, $Y' \leq Y$ (i.e.\ $\mathrm{Omb}(Y) \cap \mathrm{Omb}(X') = \mathrm{Omb}\, Y_{X'}$ où $Y_{X'} \overset{\mathrm{def}}{=} \widetilde{Y} \cap \mathrm{Omb}\, X'$). De plus, si $Y \leq X$, OPS $Y \leq X'$.}\note{le schéma de Mag 3 dispose $Y \ll X$ au-dessus de $Y' \ll X'$, reliés par des $\leq$ verticaux ($X' \leq X$, $Y' \leq Y$) ; le $Y'$ du bas est écrit sur un autre signe}

\struck{\emph{Mag 4}}

\marginal{\struck{Variante (optimiste) : si $Z \mathrel{\overset{\circ}{\ll}} Y$, $Z \ll X'$, alors $Y \ll X'$, i.e.\ \ill{} relié ; alors $Z$ \ill{} \ill{}}}\note{note écrite en biais dans la marge gauche, en face de Mag 3, barrée avec le reste}
\marginal{i.e.\ $\exists\, \ill{}$ \ill{} $Z$ dans $Y$}\note{courte note en biais dans la marge gauche, en face de Mag 2}

\page{110}
\emph{Mag 1} $X \leq Y$ et $X \ll Y$ sont des relations d'ordre, et
\[
  X \leq Y \Longrightarrow X \ll Y
\]

\emph{Mag 2} Soient $Z \ll X$, alors l'ensemble
\[
  \widetilde{X}^{Z} = \lbrace X' \in \widetilde{X} \mid Z \ll X' \rbrace
\]
a un \emph{plus petit élément}. Si \uncertain{c'est} $X$, on écrit $Z \mathrel{\overset{\circ}{\ll}} X$.\note{après « $X'$ », un signe lourdement biffé ; le $Z$ de « $Z \ll X'$ » est surchargé}

\emph{Mag 3} $X \mathrel{|{\circ}|} Y$ est symétrique et antiréflexive, i.e.\ elle implique $Y \mathrel{|{\circ}|} X$ et $X \neq Y$. De plus on a
\[
  X \mathrel{|{\circ}|} Y,\ X' \mathrel{\overset{\circ}{\ll}} X,\ Y' \mathrel{\overset{\circ}{\ll}} Y \Longrightarrow X' \mathrel{|{\circ}|} Y'
\]

\emph{Mag 4} Si $X \in \mathcal{M}$, $Y, Z \leq X$, \struck{\ill{}} \uncertain{alors}
\[
  Y \neq Z \Longrightarrow Y \mathrel{|{\circ}|} Z
\]
\note{Mag 1 à Mag 4 reprennent, dans l'ordre 1, 2, 5'', 4'', les axiomes Mag 1 et Mag 2 barrés de la page 109 et Mag 5'', Mag 4'' des pages 102-103 ; « $Y \neq Z$ » est écrit sur un premier signe biffé}

En termes de \uncertain{ces} données, on définit
\[
  X \between Y \overset{\mathrm{def}}{\Longleftrightarrow} \forall\, X' \in \widetilde{X} \smallsetminus \widetilde{X} \cap \widetilde{Y},\ Y' \in \widetilde{Y} \smallsetminus \widetilde{X} \cap \widetilde{Y}, \text{ on a } X' \mathrel{|{\circ}|} Y'
\]
\[
  X \parallel Y \overset{\mathrm{def}}{\Longleftrightarrow} \forall\, X' \in \widetilde{X},\ Y' \in \widetilde{Y}, \text{ on a } X' \mathrel{|{\circ}|} Y'
\]
On \struck{\uncertain{a}} a alors
\[
  X \parallel Y \Longleftrightarrow X \between Y \text{ et } \widetilde{X} \cap \widetilde{Y} = \emptyset \quad (\text{i.e.\ } \lbrace X, Y \rbrace \text{ non minoré dans } \mathcal{M})
\]

\struck{\ill{}} \emph{Figure} \add{$F$} (d'un magasin) : partie $\leq$-fermée, telle que $\forall\, X, Y \in F$, on ait $X \between Y$. \emph{\struck{Atelier : Magasin}}, \emph{Ex} : les $\widetilde{X}$
\note{« Atelier : Magasin » est souligné puis barré d'un trait horizontal ; la page s'arrête sur « Ex : les $\widetilde{X}$ »}

\page{111}
\emph{Atelier} \add{(spéciaux)} : est donné par $(\mathfrak{F}, \leq, \ll ; \mathrel{|{\circ}|})$ où $\mathfrak{F}$ est un ensemble muni de deux relations d'ordre $\leq$, $\ll$, et d'une \struck{\ill{} \ill{}} relation. \struck{\ill{}} Axiomes

\emph{At} 1 $\mathfrak{F}$ stable par $\mathrm{Sup}^{\leq}$ de familles \emph{majorées} \add{$\neq \emptyset$ et $\mathfrak{F}$,} \uncertain{quelconques}.\note{« $\neq \emptyset$ et $\mathfrak{F}$, » est écrit au-dessus de la ligne et relié à « stable » ; « \struck{de} $\mathfrak{F}$ » est biffé devant « majorées », souligné}

\emph{At} 2 Toute \emph{figure} $F$ est le $\mathrm{Sup}^{\leq}$ de \add{l'ens.\ $\widetilde{F}$ de} ses « strates »
\[
  \widetilde{F} = \lbrace X \in \mathfrak{F} \mid X \leq F,\ X \text{ « strictement irréductible »} \rbrace
\]
\marginal{\ill{} $F \leq G \Longleftrightarrow \widetilde{F} \ill{}$ (condition sur les figures par strates) $\widetilde{F}$, \struck{\ill{}} \ill{} en $\leq$, \ill{} les \struck{\ill{}} \ill{} de la figure (\ill{} || \ill{} pas « \uncertain{développement} »).}\note{note écrite en biais dans la marge gauche, en face de At 2 et de la suite, lue par fragments ; un mot y est souligné de trois traits}

Si $\mathcal{M}$ désigne l'ens.\ des figures \uncertain{élémentaires} (ou plutôt str.\ irréd.\ de $\mathfrak{F}$), \struck{\ill{}} munie de la structure $\leq$ induite, \uncertain{donc} la connaissance de $(\mathfrak{F}, \leq)$ équivaut à celle de $(\mathcal{M}, \leq)$, avec un ens.\ $\mathfrak{F}$ \add{de parties fermées de $\mathcal{M}$} (les $\widetilde{F}$), contenant avec chaque \ill{} \struck{\ill{}} \ill{} partie fermée, \ill{} recouvrant $\mathcal{M}$ \ill{}, contenant au moins la partie vide de $\mathcal{M}$ (figure vide). Ceci posé :\note{les quatre lignes qui suivent At 2 sont écrites vite, avec des ajouts interlinéaires ; la construction de la phrase n'est pas établie}

\marginal{\ill{} $\mathcal{L}$ \ill{} \ill{} \ill{} $C$}\note{quelques signes en biais dans la marge gauche, en face de « $(\mathcal{M}, \leq)$ »}

\emph{At} 3 \struck{$F \between G \Longleftarrow \forall\, X \in \widetilde{F}$, $\exists\, Y \in \widetilde{G}$ avec $X \ll Y$} \add{et si $\forall\, X \in \widetilde{F}$, $Y \in \widetilde{G}$,} \struck{$X \cup Y$} $X \cup Y$ existe dans $\mathfrak{F}$ (i.e.\ $\lbrace X, Y \rbrace$ majoré pour $\leq$), alors $F \cup G$ existe dans $\mathfrak{F}$.\note{At 3 est barré d'un long trait oblique ; la première formule est de plus barrée d'un trait horizontal, et « et si $\forall\, X \in \widetilde{F}$, $Y \in \widetilde{G}$, » est écrit au-dessus. Un trait vertical dans la marge marque l'énoncé}

\struck{On introduit \uncertain{la} relation dans $\mathfrak{F}$ : $F \between G \Longrightarrow \lbrace F, G \rbrace$ majoré en $\mathfrak{F}$. Par At 3, il suffit de la connaître dans $\mathcal{M}$.}\note{sous $\lbrace F, G \rbrace$, « i.e.\ $F \cup G$ existe » ; ce passage est barré du même trait oblique}

\page{112}
\marginal{Critère de \uncertain{vérification} \ill{} \ill{}}\note{note en biais dans la marge gauche, en face de At 4, entourée d'un trait courbe}
\emph{At} 4 $F \ll G \Longleftrightarrow \forall\, X \in \widetilde{F}$, $\exists\, Y \in \widetilde{G}$, avec $X \ll Y$

Donc la connaissance de $\ll$ sur $\mathfrak{F}$ équivaut à celle de $\ll$ sur $\mathcal{M}$.

\struck{\emph{At} 5 $F \mathrel{|{\circ}|} G$ implique que $F, G \in \mathcal{M}$.}\note{devant « At 5 », un premier « At » lourdement biffé}

\marginal{Critère de compatibilité \uncertain{par les strates}}\note{note en biais dans la marge gauche, en face de At 3, reliée par une grande accolade à la ligne « + les axiomes Mag 1 -- Mag 4 »}
\emph{At} 3 \add{(spéc.)} (a) Soient $F, G \in \mathfrak{F}$. Pour que $\lbrace F, G \rbrace$ majoré dans $\mathfrak{F}^{\leq}$ i.e.\ $F \cup G$ existe, il faut et il suffit que \struck{$\forall\, X \in \widetilde{F}$, $Y \in \widetilde{G}$}
\[
  \forall\, X \in \widetilde{F} \smallsetminus \widetilde{F} \cap \widetilde{G},\ Y \in \widetilde{G} \smallsetminus \widetilde{F} \cap \widetilde{G} \Longrightarrow X \mathrel{|{\circ}|} Y
\]
\struck{D'autre part, la} (b) \struck{si} $F \mathrel{|{\circ}|} G \Longrightarrow F, G \in \mathcal{M}$, i.e.\ $\mathrel{|{\circ}|}$ est en fait une relation dans $\mathcal{M}$.\note{les lettres (a), (b) sont cerclées dans le manuscrit. (b) reprend l'At 5 biffé plus haut ; « D'autre part, la » et « si » sont barrés}

+ les axiomes Mag 1 -- Mag 4 sur $(\mathcal{M}, \leq, \ll, \mathrel{|{\circ}|})$.

En termes de cette dernière structure, \add{associée au magasin} la donnée d'un \emph{atelier} (\uncertain{spécial}) : la donnée de $\mathfrak{F} \subset \mathfrak{P}(\mathcal{M})$, satisfaisant les conditions

\emph{At $\mathcal{M}$ 1} $\forall\, \Phi \in \mathfrak{F}$, $\Phi$ est fermé pour $\leq$

\emph{At $\mathcal{M}$ 2} Si $\Phi \in \mathfrak{F}$, $\Psi \subset \Phi$, \struck{alors} est une partie fermée de $\Phi$, alors $\Psi \in \mathfrak{F}$

\emph{At $\mathcal{M}$ 3} (a) $\bigcup_{\Phi \in \mathfrak{F}} \Phi = \mathcal{M}$, et (b) $\mathfrak{F} \neq \emptyset$ (ce dernier \uncertain{automatique} si $\mathcal{M} \neq \emptyset$), ou encore (b') $\emptyset \in \mathfrak{F}$.\note{(a), (b), (b') cerclés dans le manuscrit ; « est une » est écrit au-dessus de « alors », biffé, dans At $\mathcal{M}$ 2}

\page{113}
\emph{At $\mathcal{M}$ 4} Si $\Phi, \Psi \in \mathfrak{F}$, alors $\Phi \cup \Psi \in \mathfrak{F}$ ssi \struck{$\forall\, X \in \Phi$, $Y \in \Psi$, \ill{} $X, Y \notin \Phi \cap \Psi$, on a $X \mathrel{|{\circ}|} Y$} \add{$X \between Y$, i.e.} $\forall\, X' \in \widetilde{X} \smallsetminus \widetilde{X} \cap \widetilde{Y}$, $Y' \in \widetilde{Y} \smallsetminus \widetilde{X} \cap \widetilde{Y}$, \ill{} \struck{\ill{}} $X' \mathrel{|{\circ}|} Y'$.\note{la ligne « $\forall\, X \in \Phi$, $Y \in \Psi$, \ldots, on a $X \mathrel{|{\circ}|} Y$ » est barrée ; au-dessus, « $X \between Y$, i.e. » et, au-dessous, la condition sur $X'$, $Y'$ ; les variables $X$, $Y$ qui y restent, sans quantificateur, sont sur la page}

L'atelier est dit \struck{fini} « à figures finies » si les multiplicités de ses figures sont finies. Cela revient à se donner \struck{\ill{}} le magasin correspondant, et prendre alors
\[
  \mathfrak{F} = \lbrace \Phi \in \mathfrak{P}_f(\mathcal{M}) \mid X, Y \in \Phi \Rightarrow X \between Y \rbrace ,
\]
mais \struck{\ill{}} cela satisfait aux conditions d'un \uncertain{M-atelier} \ill{} \ill{}.\note{« $\Phi$ » est chaque fois écrit sur une autre lettre biffée}

\emph{\struck{At $\mathcal{M}$} Mag fin} $\forall\, X \in \mathcal{M}$, $\widetilde{X}$ est fini. \struck{\uncertain{(magasin fini, multistrates finies)}}\note{la parenthèse finale, deux lignes serrées, est barrée de deux traits}

(Axiome des \emph{magasins} \add{dans cette partie algébrique,} au cours du travail), \add{\uncertain{c'est} finalement \emph{l'atelier}} qui est au centre de la structure constamment, alors que l'ensemble $\mathfrak{F}$ de toutes les figures joue un rôle un tantinet accessoire.\note{la page, et le dossier, s'arrêtent ici, au premier tiers de la feuille}

\end{document}
